\documentclass[10pt,a4paper]{amsart}
\usepackage{amsmath,amssymb,amsthm,graphicx}
% Retained minimal preamble for the Stacks Project chapter more-morphisms.tex (More on Morphisms).
% Verbatim excerpt of the chapter preamble: Xy-pic setup (native preamble lines 14-18) and the native
% theorem declarations with the mathematical macros (native preamble lines 154-217). Copied mechanically;
% nothing added, removed or redefined. The wrapper class and the amsmath/amssymb/amsthm/graphicx base
% packages are supplied by the assembler, so the class/externaldocument part of the native preamble is omitted.
% For commutative diagrams we use Xy-pic
\usepackage[all]{xy}

% We use 2cell for 2-commutative diagrams.
\xyoption{2cell}
% Theorem environments.
%
\theoremstyle{plain}
\newtheorem{theorem}[subsection]{Theorem}
\newtheorem{proposition}[subsection]{Proposition}
\newtheorem{lemma}[subsection]{Lemma}

\theoremstyle{definition}
\newtheorem{definition}[subsection]{Definition}
\newtheorem{example}[subsection]{Example}
\newtheorem{exercise}[subsection]{Exercise}
\newtheorem{situation}[subsection]{Situation}

\theoremstyle{remark}
\newtheorem{remark}[subsection]{Remark}
\newtheorem{remarks}[subsection]{Remarks}

\numberwithin{equation}{subsection}

% Macros
%
\def\lim{\mathop{\mathrm{lim}}\nolimits}
\def\colim{\mathop{\mathrm{colim}}\nolimits}
\def\Spec{\mathop{\mathrm{Spec}}}
\def\Hom{\mathop{\mathrm{Hom}}\nolimits}
\def\Ext{\mathop{\mathrm{Ext}}\nolimits}
\def\SheafHom{\mathop{\mathcal{H}\!\mathit{om}}\nolimits}
\def\SheafExt{\mathop{\mathcal{E}\!\mathit{xt}}\nolimits}
\def\Sch{\mathit{Sch}}
\def\Mor{\mathop{\mathrm{Mor}}\nolimits}
\def\Ob{\mathop{\mathrm{Ob}}\nolimits}
\def\Sh{\mathop{\mathit{Sh}}\nolimits}
\def\NL{\mathop{N\!L}\nolimits}
\def\CH{\mathop{\mathrm{CH}}\nolimits}
\def\proetale{{pro\text{-}\acute{e}tale}}
\def\etale{{\acute{e}tale}}
\def\QCoh{\mathit{QCoh}}
\def\Ker{\mathop{\mathrm{Ker}}}
\def\Im{\mathop{\mathrm{Im}}}
\def\Coker{\mathop{\mathrm{Coker}}}
\def\Coim{\mathop{\mathrm{Coim}}}

% Boxtimes
%
\DeclareMathSymbol{\boxtimes}{\mathbin}{AMSa}{"02}

%
% Macros for moduli stacks/spaces
%
\def\QCohstack{\mathcal{QC}\!\mathit{oh}}
\def\Cohstack{\mathcal{C}\!\mathit{oh}}
\def\Spacesstack{\mathcal{S}\!\mathit{paces}}
\def\Quotfunctor{\mathrm{Quot}}
\def\Hilbfunctor{\mathrm{Hilb}}
\def\Curvesstack{\mathcal{C}\!\mathit{urves}}
\def\Polarizedstack{\mathcal{P}\!\mathit{olarized}}
\def\Complexesstack{\mathcal{C}\!\mathit{omplexes}}
% \Pic is the operator that assigns to X its picard group, usage \Pic(X)
% \Picardstack_{X/B} denotes the Picard stack of X over B
% \Picardfunctor_{X/B} denotes the Picard functor of X over B
\def\Pic{\mathop{\mathrm{Pic}}\nolimits}
\def\Picardstack{\mathcal{P}\!\mathit{ic}}
\def\Picardfunctor{\mathrm{Pic}}
\def\Deformationcategory{\mathcal{D}\!\mathit{ef}}
\usepackage{cleveref}
\crefname{originallemma}{original lemma}{original lemmas}
\Crefname{originallemma}{Original lemma}{Original lemmas}
\makeatletter
\newlabel{lemma-H1-O-picard-flat-proper}{{lemma-H1-O-picard-flat-proper}{source}{}{sourcepointer-1}{}}
\newlabel{lemma-H1-O-picard-flat-proper@cref}{{[originallemma][0][]lemma-H1-O-picard-flat-proper}{[1][0][]source}{}{}{}}
\makeatother
\crefname{originallemma}{original lemma}{original lemmas}
\Crefname{originallemma}{Original lemma}{Original lemmas}
\makeatletter
\newlabel{lemma-diagonal-picard-flat-proper}{{lemma-diagonal-picard-flat-proper}{source}{}{sourcepointer-2}{}}
\newlabel{lemma-diagonal-picard-flat-proper@cref}{{[originallemma][0][]lemma-diagonal-picard-flat-proper}{[1][0][]source}{}{}{}}
\makeatother
\crefname{originallemma}{original lemma}{original lemmas}
\Crefname{originallemma}{Original lemma}{Original lemmas}
\makeatletter
\newlabel{more-algebra-lemma-lift-perfect-from-residue-field}{{more-algebra-lemma-lift-perfect-from-residue-field}{source}{}{sourcepointer-3}{}}
\newlabel{more-algebra-lemma-lift-perfect-from-residue-field@cref}{{[originallemma][0][]more-algebra-lemma-lift-perfect-from-residue-field}{[1][0][]source}{}{}{}}
\makeatother
\crefname{originallemma}{original lemma}{original lemmas}
\Crefname{originallemma}{Original lemma}{Original lemmas}
\makeatletter
\newlabel{perfect-lemma-affine-compare-bounded}{{perfect-lemma-affine-compare-bounded}{source}{}{sourcepointer-4}{}}
\newlabel{perfect-lemma-affine-compare-bounded@cref}{{[originallemma][0][]perfect-lemma-affine-compare-bounded}{[1][0][]source}{}{}{}}
\makeatother
\crefname{originallemma}{original lemma}{original lemmas}
\Crefname{originallemma}{Original lemma}{Original lemmas}
\makeatletter
\newlabel{perfect-lemma-flat-proper-perfect-direct-image-general}{{perfect-lemma-flat-proper-perfect-direct-image-general}{source}{}{sourcepointer-5}{}}
\newlabel{perfect-lemma-flat-proper-perfect-direct-image-general@cref}{{[originallemma][0][]perfect-lemma-flat-proper-perfect-direct-image-general}{[1][0][]source}{}{}{}}
\makeatother
\crefname{originallemma}{original lemma}{original lemmas}
\Crefname{originallemma}{Original lemma}{Original lemmas}
\makeatletter
\newlabel{perfect-lemma-perfect-affine}{{perfect-lemma-perfect-affine}{source}{}{sourcepointer-6}{}}
\newlabel{perfect-lemma-perfect-affine@cref}{{[originallemma][0][]perfect-lemma-perfect-affine}{[1][0][]source}{}{}{}}
\makeatother
\title{Vanishing of the first cohomology of multiples of a relatively ample bundle on a flat proper family}
\author{The Stacks Project Authors}
\date{stacks-project:more-morphisms.tex:local-snapshot-20261011; GNU FDL 1.2 or later}
\begin{document}
\maketitle
\noindent\textit{Editorial note.} Statement span is lines 9518-9541 of sources/stacks-more-morphisms/pinned-original-source.tex; proof span is lines 9543-9619, the complete author proof (75 body lines).\par\medskip
\noindent\textit{Editorial note.} Navigation stubs: 6 cited-but-not-included labels are supplied as non-load-bearing pointers with their pinned tags where the snapshot tag index has them.\par\medskip
\noindent\textit{Editorial note.} extra\_packages is empty on purpose: the retained preamble already loads xy with options and a second declaration clashes in the assembler wrapper.\par\medskip
\noindent\textit{Editorial note.} Text is retained verbatim, including the author's own typos and placeholders.\par\medskip
\noindent\textit{Allowed standard prerequisites (editorial).}\begin{itemize}
\item Explicit prerequisite boundary. Relative ampleness, cohomology and base change for flat proper morphisms, and finite local freeness criteria; the uniform-threshold argument is retained.
\end{itemize}\medskip
\ifdefined\hypertarget\hypertarget{sourcepointer-1}{}\fi
\noindent\textit{Source reference.} Original lemma lemma-H1-O-picard-flat-proper: sources/stacks-more-morphisms/pinned-original-source.tex, source lines 9430--9451. Cited by the retained statement or proof but not part of the collected spans; supplied only as a navigation pointer. The argument does not depend on its content.\par\medskip
\ifdefined\hypertarget\hypertarget{sourcepointer-2}{}\fi
\noindent\textit{Source reference.} Original lemma lemma-diagonal-picard-flat-proper: sources/stacks-more-morphisms/pinned-original-source.tex, source lines 9199--9215. Cited by the retained statement or proof but not part of the collected spans; supplied only as a navigation pointer. The argument does not depend on its content.\par\medskip
\ifdefined\hypertarget\hypertarget{sourcepointer-3}{}\fi
\noindent\textit{Source reference.} Original lemma more-algebra-lemma-lift-perfect-from-residue-field: sources/stacks-more-morphisms/pinned-original-source.tex, source lines 9619--9619. Cited by the retained statement or proof but not part of the collected spans; supplied only as a navigation pointer. The argument does not depend on its content.\par\medskip
\ifdefined\hypertarget\hypertarget{sourcepointer-4}{}\fi
\noindent\textit{Source reference.} Original lemma perfect-lemma-affine-compare-bounded: sources/stacks-more-morphisms/pinned-original-source.tex, source lines 9619--9619. Cited by the retained statement or proof but not part of the collected spans; supplied only as a navigation pointer. The argument does not depend on its content.\par\medskip
\ifdefined\hypertarget\hypertarget{sourcepointer-5}{}\fi
\noindent\textit{Source reference.} Original lemma perfect-lemma-flat-proper-perfect-direct-image-general: sources/stacks-more-morphisms/pinned-original-source.tex, source lines 9619--9619. Cited by the retained statement or proof but not part of the collected spans; supplied only as a navigation pointer. The argument does not depend on its content.\par\medskip
\ifdefined\hypertarget\hypertarget{sourcepointer-6}{}\fi
\noindent\textit{Source reference.} Original lemma perfect-lemma-perfect-affine: sources/stacks-more-morphisms/pinned-original-source.tex, source lines 9619--9619. Cited by the retained statement or proof but not part of the collected spans; supplied only as a navigation pointer. The argument does not depend on its content.\par\medskip
\begin{lemma}
\label{lemma-H1-O-multiple-picard-flat-proper}
Consider $n$ commutative diagrams of schemes
$$
\xymatrix{
X_i \ar[rr] \ar[dr]_{f_i} & & X \ar[dl]^f \\
& S
}
$$
with $f_i : X_i \to S$ and $f : X \to S$ satisfying the hypotheses of
Lemma \ref{lemma-diagonal-picard-flat-proper}.
Let $\mathcal{L}$ be an invertible $\mathcal{O}_X$-module
and let $\mathcal{L}_i$ be the pullback to $X_i$. Let $Z \subset S$,
resp.\ $Z_i \subset S$ be the locally closed subscheme constructed
in Lemma \ref{lemma-diagonal-picard-flat-proper}
for $(f, \mathcal{L})$, resp.\ $(f_i, \mathcal{L}_i)$
so that $Z \subset \bigcap_{i = 1, \ldots, n} Z_i$. If $s \in Z$ and
$$
H^1(X_s, \mathcal{O}) \longrightarrow
\bigoplus\nolimits_{i = 1, \ldots, n} H^1(X_{i, s}, \mathcal{O})
$$
is injective, then $Z \cap U = (\bigcap_{i = 1, \ldots, n} Z_i) \cap U$
(scheme theoretic intersection) for some open neighbourhood $U$ of $s$.
\end{lemma}

\begin{proof}
This lemma is a variant of Lemma \ref{lemma-H1-O-picard-flat-proper}
and we strongly urge the reader to read that proof first; this proof
is basically a copy of that proof with minor modifications. It follows
from the description of (scheme valued) points of $Z$ and the $Z_i$
that $Z \subset \bigcap_{i = 1, \ldots, n} Z_i$ where we take the
scheme theoretic intersection. Thus we may replace $S$ by the scheme
theoretic intersection $\bigcap_{i = 1, \ldots, n} Z_i$. After shrinking
$S$ to an affine open neighbourhood of $s$ we may assume that
$\mathcal{L}_i = \mathcal{O}_{X_i}$ for $i = 1, \ldots, n$.
Let $E = Rf_*\mathcal{L}$ and
$E_i = Rf_{i, *}\mathcal{L}_i = Rf_{i, *}\mathcal{O}_{X_i}$.
These are perfect complexes whose formation commutes with arbitrary
change of base (Derived Categories of Schemes, Lemma
\ref{perfect-lemma-flat-proper-perfect-direct-image-general}).
In particular we see that
$$
E \otimes_{\mathcal{O}_S}^\mathbf{L} \kappa(s) =
R\Gamma(X_s, \mathcal{L}_s) = R\Gamma(X_s, \mathcal{O}_{X_s})
$$
The second equality because $s \in Z$. Set
$h_j = \dim_{\kappa(s)} H^j(X_s, \mathcal{O}_{X_s})$.
After shrinking $S$ we can represent $E$ by a complex
$$
\mathcal{O}_S \to \mathcal{O}_S^{\oplus h_1} \to
\mathcal{O}_S^{\oplus h_2} \to \ldots
$$
see More on Algebra, Lemma
\ref{more-algebra-lemma-lift-perfect-from-residue-field}
(strictly speaking this also uses
Derived Categories of Schemes, Lemmas
\ref{perfect-lemma-affine-compare-bounded} and
\ref{perfect-lemma-perfect-affine}). Similarly, we may assume $E_i$
is represented by a complex
$$
\mathcal{O}_S \to \mathcal{O}_S^{\oplus h_{i, 1}} \to
\mathcal{O}_S^{\oplus h_{i, 2}} \to \ldots
$$
where $h_{i, j} = \dim_{\kappa(s)} H^j(X_{i, s}, \mathcal{O}_{X_{i, s}})$.
By functoriality of cohomology we have a map
$$
E \longrightarrow E_i
$$
in $D(\mathcal{O}_S)$ whose formation commutes with change of base.
Since the complex representing $E$ is a finite complex of finite free
modules and since $S$ is affine, we can choose a map of complexes
$$
\xymatrix{
\mathcal{O}_S \ar[r]_d \ar[d]_{a_i} &
\mathcal{O}_S^{\oplus h_1} \ar[r] \ar[d]_{b_i} &
\mathcal{O}_S^{\oplus h_2} \ar[r] \ar[d]_{c_i} & \ldots \\
\mathcal{O}_S \ar[r]^{d_i} &
\mathcal{O}_S^{\oplus h_{i, 1}} \ar[r] &
\mathcal{O}_S^{\oplus h_{i, 2}} \ar[r] & \ldots
}
$$
representing the given map $E \to E_i$. Since $s \in Z$ we see that
the trivializing section of $\mathcal{L}_s$ pulls back to a trivializing
section of $\mathcal{L}_{i, s} = \mathcal{O}_{X_{i, s}}$. Thus
$a_i \otimes \kappa(s)$ is an isomorphism, hence after shrinking $S$
we see that $a_i$ is an isomorphism. Finally, we use the hypothesis
that $H^1(X_s, \mathcal{O}) \to
\bigoplus_{i = 1, \ldots, n} H^1(X_{i, s}, \mathcal{O})$
is injective, to see that there exists a $h_1 \times h_1$ minor of the
matrix defining $\oplus b_i$ which maps to a nonzero
element in $\kappa(s)$. Hence after shrinking $S$ we may assume that
$(b_1, \ldots, b_n) : \mathcal{O}_S^{h_1}
\to \bigoplus_{i = 1, \ldots, n} \mathcal{O}_S^{h_{i, 1}}$
is injective. However, since $\mathcal{L}_i = \mathcal{O}_{X_i}$
we see that $d_i = 0$ for $i = 1, \ldots n$. It follows that $d = 0$
because $(b_1, \ldots, b_n) \circ d = (\oplus d_i) \circ (a_1, \ldots, a_n)$.
In this way we see
that the trivializing section of $\mathcal{L}_s$ lifts to a section
of $\mathcal{L}$ over $X$. A straightforward topological argument (omitted)
shows that this means that $\mathcal{L}$ is trivial after possibly
shrinking $S$ a bit further.
\end{proof}

\end{document}
